\documentclass{beamer}

\mode<presentation>
{
  \usetheme{Warsaw}
  \setbeamercovered{transparent}
}

\usetheme{Warsaw}
\setbeamersize{text margin left=4mm}
\setbeamersize{text margin right=4mm}

\usepackage[utf8]{inputenc}
\usepackage[backend=bibtex,style=authoryear]{biblatex}
\addbibresource{thesis.bib}

\usepackage{graphicx}
\usepackage{booktabs}
\usepackage{adjustbox}
\usepackage{capt-of}
\usepackage{threeparttable}
\usepackage{amsmath}
\usepackage{amssymb}
\usepackage{xcolor}
\usepackage{tikz}
\usepackage{appendixnumberbeamer}
\usepackage{dirtytalk}

\title[Dissertation Defense]{Essays on Heterogeneous Returns, Trust, and Household Wealth}
\author[Edwards]{Decory Edwards}
\institute[JHU]{Johns Hopkins University}
\date[June 22, 2026]{Final Dissertation Defense \\ June 22, 2026}

\newcommand{\assetpath}[1]{assets/#1}
\newcommand{\floatingnavbutton}[2]{%
  \begin{tikzpicture}[remember picture,overlay]
    \node[anchor=south east] at ([xshift=-0.5em,yshift=0.8em]current page.south east) {\hyperlink{#1}{\beamerbutton{#2}}};
  \end{tikzpicture}%
}
\newcommand{\centernavbutton}[2]{%
  \begin{center}
    \hyperlink{#1}{\beamerbutton{#2}}
  \end{center}
}
\newcommand{\slidetable}[1]{%
  \begin{center}
  \begin{adjustbox}{max totalsize={\textwidth}{0.70\textheight}}
  #1
  \end{adjustbox}
  \end{center}
}


\begin{document}

\begin{frame}
  \titlepage
\end{frame}

\begin{frame}{Roadmap}
  \tableofcontents[hideallsubsections]
\end{frame}

\section{Introduction}

\begin{frame}{Introduction}
  \begin{enumerate}
    \item \textbf{Chapter 1}: Can return heterogeneity help a standard heterogeneous-agent model match empirical wealth moments?
    \item \textbf{Chapter 2}: Is there a systematic relationship between trust and returns to net wealth, and can we still estimate a persistent household return component?
    \item \textbf{Chapter 3}: If wealth is tied to social institutions and historical group structure, does the standard objective benchmark understate wealth inequality?
  \end{enumerate}

  \vspace{0.6em}
  \textit{Taken together, the chapters move from explaining wealth concentration, to explaining realized performance, to measuring inequality itself.}
\end{frame}

\section{Chapter 1}

\begin{frame}{Chapter 1: Matching Wealth Moments with Heterogeneous Returns}
  \begin{itemize}
    \item I incorporate return heterogeneity into a standard heterogeneous-agent model and estimate the return distribution needed to match the wealth distribution.
    \item Labor income uncertainty and life-cycle structure explain part of wealth dispersion, but not enough to reproduce the empirical skewness of wealth.
    \item Adding return heterogeneity yields a much better fit to SCF wealth moments and implies a more realistic aggregate MPC.
  \end{itemize}

  \vspace{0.6em}
  Empirically, this is motivated by evidence that realized household returns vary widely across people \parencite{aflgdmlp20, Daminato2024}.
\end{frame}

\begin{frame}[label=ch1-results]{Chapter 1: Main Result}
  \begin{columns}
    \column{0.5\textwidth}
    \centering
    \includegraphics[width=\textwidth,height=0.42\textheight,keepaspectratio]{\assetpath{ch1_lorenz_point.png}}
    \column{0.5\textwidth}
    \centering
    \includegraphics[width=\textwidth,height=0.42\textheight,keepaspectratio]{\assetpath{ch1_lorenz_dist.png}}
  \end{columns}

  \vspace{0.3em}
  \scriptsize
  \begin{center}
  \begin{tabular}{lcccccc}
    \toprule
     & \textbf{0--20} & \textbf{20--40} & \textbf{40--60} & \textbf{60--80} & \textbf{80--95} & \textbf{Top 5\%} \\
    \midrule
    Data             & $-0.002$ & 0.011 & 0.044 & 0.118 & 0.255 & 0.574 \\
    No heterogeneity & 0.030    & 0.085 & 0.142 & 0.224 & 0.285 & 0.233 \\
    Heterogeneity    & 0.005    & 0.019 & 0.042 & 0.096 & 0.251 & 0.588 \\
    \bottomrule
  \end{tabular}
  \end{center}
  \vspace{0.2em}
  \tiny\textit{Wealth shares by net-worth percentile, infinite-horizon (PY) model. Estimated uniform return distribution: $R=1.0204$, $\nabla=0.06833$.}
\end{frame}

\begin{frame}[label=ch1-policy]{Chapter 1: Policy Implications — Wealth Distribution}
  \begin{columns}
    \column{0.55\textwidth}
    \begin{center}
      \includegraphics[width=\textwidth,height=0.68\textheight,keepaspectratio]{\assetpath{ch1_lorenz_tax_py.png}}
    \end{center}

    \column{0.45\textwidth}
    \scriptsize
    \begin{itemize}
      \item Two revenue-equivalent policies (each raising 1\% of GDP): wealth tax $1+r-\tau_w$ lowers the gross return for all; capital income tax $1+(1-\tau_{ci})r$ taxes only positive returns.
      \item The CIT reduces wealth concentration by more than the WT in both models: only households with strictly positive returns face the capital income tax, a smaller tax pool than the wealth tax.
    \end{itemize}
  \end{columns}
\end{frame}

\begin{frame}[label=ch1-welfare]{Chapter 1: Policy Implications — Welfare}
  \scriptsize
  \begin{itemize}
    \item Welfare is measured by consumption-equivalent $\Delta$: the \% increase in consumption under the WT that makes a newborn --- before drawing their return type --- indifferent to the CIT. $\Delta > 0$ means CIT preferred.
    \item \textbf{Aggregate:} $\Delta = 0.20\%$ (PY) and $0.15\%$ (LC) --- the CIT is preferred on average for a newborn.
    \item \textbf{By return type:} 6 of 7 types prefer the CIT; only the highest-return type ($R=1.079$) prefers the WT. Low-return types gain most from switching to CIT because the WT reduces their already-low gross return while generating little revenue from them.
  \end{itemize}
  \vspace{0.3em}
  \slidetable{
  \begin{tabular}{lccccccc}
    \toprule
    & \textbf{Type 1} & \textbf{Type 2} & \textbf{Type 3} & \textbf{Type 4} & \textbf{Type 5} & \textbf{Type 6} & \textbf{Type 7} \\
    \midrule
    Baseline $R$ & 0.962 & 0.981 & 1.001 & 1.020 & 1.040 & 1.059 & 1.079 \\
    CE $\Delta$ (WT vs CIT) & +0.23\% & +0.27\% & +0.34\% & +0.32\% & +0.30\% & +0.24\% & $-$0.31\% \\
    \bottomrule
  \end{tabular}
  }
\end{frame}


\section{Chapter 2}

\begin{frame}{Chapter 2: Moderate Trust and Maximum Performance on Financial Investments}
  \begin{itemize}
    \item Economic performance is hump-shaped in trust: moderate trust maximizes outcomes at the individual level.
    \item Using the RAND HRS panel --- which includes both household financial data and a generalized trust measure --- we test whether this holds for returns to net wealth.
    \item It does: the return-maximizing trust level is 6--7 on a 1--10 scale, and the panel structure allows us to estimate a persistent household return component.
  \end{itemize}
\end{frame}

\begin{frame}[label=ch2-results]{Chapter 2: Main Result}
  \scriptsize
  \begin{itemize}
    \item Trust and Trust$^2$ are jointly significant: returns to net wealth are hump-shaped in trust.
    \item The result is robust across pooled OLS and correlated random effects (CRE/Mundlak) specifications.
  \end{itemize}
  \vspace{0.2em}
  \slidetable{
  \begin{tabular}{lcccc}
    \toprule
    & \multicolumn{2}{c}{Pooled OLS} & \multicolumn{2}{c}{CRE} \\
    \cmidrule(lr){2-3}\cmidrule(lr){4-5}
    & (1) Baseline & (2) Extended & (3) Baseline & (4) Extended \\
    \midrule
    Trust     & $0.03^{***}$ & $0.03^{***}$ & $0.03^{***}$ & $0.03^{***}$ \\
              & $(0.01)$     & $(0.01)$     & $(0.01)$     & $(0.01)$     \\
    Trust$^2$ & $-0.00^{**}$ & $-0.00^{**}$ & $-0.00^{***}$& $-0.00^{***}$\\
              & $(0.00)$     & $(0.00)$     & $(0.00)$     & $(0.00)$     \\
    \midrule
    $N$            & 2{,}919 & 2{,}899 & 2{,}919 & 2{,}899 \\
    $R^2$          & 0.04    & 0.04    & 0.15    & 0.15    \\
    $\rho$         &         &         & 0.16    & 0.15    \\
    $\sigma_u$     &         &         & 0.16    & 0.15    \\
    $\sigma_e$     &         &         & 0.36    & 0.36    \\
    \bottomrule
    \multicolumn{5}{l}{\textit{Controls: age, wealth, education, year; col.\ 2 \& 4 add demographics.}}
  \end{tabular}}
  \floatingnavbutton{ch2-trustmax-deriv}{Where's the max?}
\end{frame}

\begin{frame}[label=ch2-trustmax-deriv]{Chapter 2: How is Maximal Trust Estimated?}
  \scriptsize
  Starting from the panel regression with controls held fixed,
  $$ r^{net}_{it} = \alpha_i + \lambda_t + X_{it}'\beta + \beta_1\,\textit{Trust}_i + \beta_2\,\textit{Trust}_i^2 + \varepsilon_{it}, $$
  the implied marginal effect of trust on returns is
  $$ \frac{\partial r^{net}_{it}}{\partial \textit{Trust}_i} = \beta_1 + 2\beta_2\,\textit{Trust}_i. $$

  \vspace{0.4em}
  Setting this derivative to zero gives the critical point
  $$ \textit{Trust}^{*} = -\frac{\beta_1}{2\beta_2}. $$

  \vspace{0.4em}
  When $\hat\beta_1 \geq 0$ and $\hat\beta_2 \leq 0$, the fitted relationship in trust is concave (hump-shaped), so $\textit{Trust}^*$ is a maximum: returns rise in trust up to $\textit{Trust}^*$ and fall beyond it. Plugging in the estimated coefficients gives the values reported on the previous slide (6--7 across specifications).

  \vspace{0.6em}
  \centernavbutton{ch2-results}{Back to Main Result}
\end{frame}

\begin{frame}[label=ch2-trustmax]{Chapter 2: Return-Maximizing Trust Level}
  \begin{columns}
    \column{0.55\textwidth}
    \begin{center}
      \includegraphics[width=\textwidth,height=0.65\textheight,keepaspectratio]{\assetpath{ch2_trustdist_clean.png}}
    \end{center}

    \column{0.45\textwidth}
    \scriptsize
    \slidetable{
    \begin{tabular}{lcc}
      \toprule
      Specification & No controls & With controls \\
      \midrule
      Avg.\ returns & 6.70 & 7.07 \\
      Pooled OLS    & 6.89 & 7.61 \\
      CRE           & 6.34 & 6.59 \\
      \bottomrule
    \end{tabular}}
    \vspace{0.4em}
    \begin{itemize}
      \item Estimated return-maximizing trust is consistently 6--7 across all specifications.
      \item A large mass of respondents is concentrated around the estimated optimal trust level.
    \end{itemize}
  \end{columns}
\end{frame}

\begin{frame}[label=ch2-persistent]{Chapter 2: Persistent Return Component}
  \scriptsize
  \begin{itemize}
    \item The estimated persistent component across individuals differs because FE absorbs all time-invariant variation (trust, demographics) directly into the fixed effect, while CRE estimates it separately via the Mundlak terms.
    \item A substantially larger share of total return variance loads onto the persistent component under FE ($\rho=0.65$) than under CRE ($\rho=0.15$).
    \item The individual effect's standard deviation is also far larger under FE ($\sigma_u=0.45$) than under CRE ($\sigma_u=0.15$).
  \end{itemize}
  \vspace{0.4em}
  \begin{center}
  \begin{tabular}{lcc}
    \toprule
     & FE & CRE \\
    \midrule
    $\rho$      & 0.65 & 0.15 \\
    $\sigma_u$  & 0.45 & 0.15 \\
    $\sigma_e$  & 0.33 & 0.36 \\
    \bottomrule
  \end{tabular}
  \end{center}
  \floatingnavbutton{ch2-persistent-hist}{See histograms}
\end{frame}

\begin{frame}[label=ch2-persistent-hist]{Chapter 2: Estimated Individual Effects --- FE vs.\ CRE}
  \begin{columns}
    \column{0.50\textwidth}
    \begin{center}
      \includegraphics[width=\textwidth,height=0.55\textheight,keepaspectratio]{\assetpath{ch2_fe_hist_main.png}}\\[2pt]
      \scriptsize\textit{FE: estimated individual effects (full controls).}
    \end{center}
    \column{0.50\textwidth}
    \begin{center}
      \includegraphics[width=\textwidth,height=0.55\textheight,keepaspectratio]{\assetpath{ch2_cre_hist_main.png}}\\[2pt]
      \scriptsize\textit{CRE: estimated individual effects (full controls).}
    \end{center}
  \end{columns}
  \vspace{0.6em}
  \centernavbutton{ch2-persistent}{Back}
\end{frame}

\section{Chapter 3}

\begin{frame}{Chapter 3: Measuring Wealth Inequality Under Ambiguity}
  \begin{columns}
    \column{0.55\textwidth}
    \begin{center}
      \includegraphics[width=\textwidth,height=0.58\textheight,keepaspectratio]{\assetpath{ch3_gap_motivation.png}}
    \end{center}

    \column{0.45\textwidth}
    \scriptsize
    \begin{itemize}
      \item Wealth is a complex object: it is far more skewed than income, and it requires enforceable ownership rights and legal/political institutions.
      \item Racial wealth inequality has a well-documented history in this country, and has generally risen over time.
      \item I argue that insights from the choice-under-uncertainty literature can be used to propose a measure of inequality aversion that addresses these complexities in wealth and its allocation across individuals.
    \end{itemize}
  \end{columns}
\end{frame}

\begin{frame}{Chapter 3: A Thought Experiment}
  \scriptsize
  \begin{itemize}
    \item To see this, consider two budget-balanced wealth transfers in a population of 60 (10 Black, 50 White; each group split into rich and poor):
    \begin{itemize}
      \item \textbf{Policy A:} take 20 units from each of 5 rich Black households; give 4 units to each of 25 poor White households.
      \item \textbf{Policy B:} take 4 units from each of 25 rich White households; give 20 units to each of 5 poor Black households.
    \end{itemize}
    \item Both move the same total wealth (100 units) and, apart from the racial labels, leave the same resulting wealth distribution --- under the standard (objective-uncertainty) ranking, Policies A and B are indistinguishable.
  \end{itemize}
  \vspace{0.6em}
  \textit{Yet it isn't hard to imagine these two policies would be treated very differently politically.}
\end{frame}

\begin{frame}{Chapter 3: Why Would They Be Viewed Differently?}
  \scriptsize
  {\setlength{\itemsep}{0.8em}
  \begin{itemize}
    \item \textbf{Historical context:} redistributive policies are about alleviating discrimination and reducing inequality. In the U.S., inequality in wealth along racial lines has a clear history. This is one reason why a voting public may not be indifferent between the two policies.
    \vspace{0.6em}
    \item \textbf{The birth lottery and Rawls:} Rawls argued that we operate behind a veil of ignorance --- we make rules without knowing which state or group we will be born into. This rationale motivates a maximin principle, which ranks social alternatives by their worst possible outcome for the group that ends up worst off.
    \vspace{0.6em}
    \item As we will see, this birth-lottery maximin logic --- maximizing the welfare of the worst-off state --- is precisely the mathematical structure underlying the choice-under-ambiguity model used to propose the alternative wealth-inequality measure in this chapter.
  \end{itemize}}
\end{frame}

\begin{frame}[label=ch3-setup]{Chapter 3: Inequality Measurement --- Setup}
  \scriptsize
  \textbf{Input:} An income frequency distribution $p(y)$ --- for each income level $y \in [0,\bar{y}]$, $p(y)$ is the share of the population at that income level. This is the standard object in the inequality literature; the objects being ranked are \textit{full distributions}, not individual outcomes.

  \vspace{0.5em}
  \textbf{Objective (Atkinson) measure.} A social planner with welfare function $U(y)$ (increasing, concave) ranks income distributions by expected utility:
  $$W^O = \int_0^{\bar{y}} U(y)\, p(y)\, dy.$$
  The \textit{equally distributed equivalent} income $y_{EDE}$ satisfies $U(y_{EDE}) = W^O$. The inequality index is:
  $$I^O = 1 - \frac{y_{EDE}}{\mu}, \qquad \mu = \text{mean income.}$$

  \vspace{0.2em}
  Using the CRRA welfare function $U(y) = \frac{y^{1-\epsilon}}{1-\epsilon}$ (with $\epsilon \geq 0$ governing inequality aversion) yields a mean-independent, unit-free measure $I^O \in [0,1]$.
\end{frame}

\begin{frame}[label=ch3-ambiguity-derive]{Chapter 3: The Ambiguity-Based Measure}
  \scriptsize
  \textbf{Reinterpreting wealth.} A wealth distribution is modeled as a \textit{state-contingent income distribution}: fix a state space $S = \{\text{Black},\,\text{White}\}$ --- the \say{birth lottery,} in the spirit of Rawls' veil of ignorance --- and let each act $f \in F$ assign to every state $s \in S$ an income frequency distribution, $f(s) = p(s)(y)$. The planner does not know the objective probability over states --- this is \textit{ambiguity} (Gilboa--Schmeidler 1989).

  \vspace{0.4em}
  The ambiguity-averse welfare criterion takes the worst-case weighting over beliefs $\lambda \in \Delta(S)$:
  $$W^A = \min_{\lambda \in \Delta(S)} \int_S \int_0^{\bar{y}} U(y)\,p(s)(y)\,dy\,d\lambda \;=\; \min_{s\,\in\, S}\;\mathbb{E}_{p(s)}\bigl[U(y)\bigr].$$
  The minimum puts all weight on whichever racial group has the lower within-group expected utility that year.

  \vspace{0.4em}
  The ambiguity EDE solves $U(w^A_{EDE}) = W^A$, giving the wealth inequality measure:
  $$I^A = 1 - \frac{w^A_{EDE}}{\mu}.$$

  \vspace{0.2em}
  Since $W^A \leq W^O$ by construction, $I^A \geq I^O$: the ambiguity criterion \textbf{systematically reports higher inequality} than the objective benchmark for the same wealth distribution.
\end{frame}

\begin{frame}[label=ch3-results]{Chapter 3: Ambiguity-Based Wealth Inequality}
  \begin{columns}
    \column{0.60\textwidth}
    \begin{center}
      \includegraphics[width=\textwidth,height=0.30\textheight,keepaspectratio]{\assetpath{ch3_measures_net_eps05.png}}\\[2pt]
      \includegraphics[width=\textwidth,height=0.30\textheight,keepaspectratio]{\assetpath{ch3_measures_gross_eps05.png}}
    \end{center}
    \vspace{0.2em}
    \tiny\textit{Shown for $\epsilon=0.5$ (moderate inequality aversion); $I^A \geq I^O$ holds for every $\epsilon \in \{0,0.25,0.5,1,1.5,2\}$.}

    \column{0.40\textwidth}
    \scriptsize
    \begin{itemize}
      \item The objective benchmark understates wealth inequality relative to the ambiguity-based measure for the same wealth distribution.
      \item This holds for both net wealth (top) and gross assets (bottom), and the measures track the broad time pattern of the median wealth-gap diagnostics.
    \end{itemize}
  \end{columns}
\end{frame}

\section{Conclusion}

\begin{frame}{Conclusion}
  \begin{itemize}
    \item \textbf{Chapter 1}: return heterogeneity can match key wealth moments in a structural macro model, with an estimated distribution close to empirical evidence.
    \item \textbf{Chapter 2}: returns to net wealth exhibit a hump-shaped relationship with trust, and the maximizing level of trust lies in a moderate range between 6 and 7.
    \item \textbf{Chapter 3}: the ambiguity-based wealth inequality measure implies more inequality than the standard objective benchmark for the same wealth distribution.
  \end{itemize}

  \vspace{0.6em}
  \textit{Overall conclusion: heterogeneous returns are central both to explaining household wealth inequality and to evaluating it.}
\end{frame}

\end{document}
